% GENERATED FILE. Edit canonical lessons, then run npm run generate:books.
% canonical-source-hash: fnv1a64:dbe564fb64600469
% canonical-lessons: LA-C161-nihil, LA-C161-nemo, LA-C161-nullus, LA-C161-uter, LA-C161-neque

\chapter{Nothing, Nobody, No, Which of Two, Neither … Nor}
\label{ch:la-nothing-nobody-no-which-of-two-neither-nor}
\hlchaptermodality{161}

\begin{chapteropening}
\textbf{By the end of this chapter:} \emph{I can say that there is nothing, nobody and none, and ask which of two.}

Complete the last lesson of chapter 161: I can say that there is nothing, nobody and none, and ask which of two.
\end{chapteropening}

\section[nihil]{nihil — nothing}

\label{lesson:LA-C161-nihil}



\begin{quote}
Before the new one: say the Latin for to arrive, then the Latin for to live (somewhere).
\end{quote}

\subsection*{You'll want to know: nihil}
\textbf{nihil} — \textquotedblleft{}nothing\textquotedblright{}. \textbf{Nihil videō.} — \textquotedblleft{}I see nothing.\textquotedblright{} Latin needs no second negative: \textbf{nihil} already means \textquotedblleft{}not anything\textquotedblright{}. English keeps it in \emph{nil} and \emph{annihilate}.

\begin{grammarlens}[title={What you've built}]
One more word for this chapter.
\end{grammarlens}

\subsection*{Your turn}
\begin{itemize}
\raggedright
  \item \emph{Say it:} \emph{nihil}
  \item \emph{Say it:} \emph{nihil}, once more
  \item \emph{Say it:} say \emph{nihil}
  \item \emph{Recall:} say the Latin for to arrive, then the Latin for to live (somewhere), then say \emph{nihil} again
\end{itemize}

\begin{tcolorbox}[breakable,colback=teal!4,colframe=teal!35!black,title={Before you move on}]
What does nihil mean? (\textquotedblleft{}Nothing\textquotedblright{}.) Say it once more.
\end{tcolorbox}
\section[nēmō, nēminis]{nēmō, nēminis — nobody, no one}

\label{lesson:LA-C161-nemo}



\begin{quote}
Before the new one: say the Latin for to live (somewhere), then the Latin for nothing.
\end{quote}

\subsection*{You'll want to know: nēmō, nēminis}
\textbf{nēmō, nēminis} — \textquotedblleft{}nobody, no one\textquotedblright{}. \textbf{Nēmō venit.} — \textquotedblleft{}Nobody is coming.\textquotedblright{} It is \textbf{ne-}, \textquotedblleft{}not\textquotedblright{}, with an old form of \textbf{homō}, \textquotedblleft{}a person\textquotedblright{}: \textquotedblleft{}not a person\textquotedblright{}.

\begin{grammarlens}[title={What you've built}]
One more word for this chapter.
\end{grammarlens}

\subsection*{Your turn}
\begin{itemize}
\raggedright
  \item \emph{Say it:} \emph{nēmō, nēminis}
  \item \emph{Say it:} \emph{nēmō, nēminis}, once more
  \item \emph{Say it:} say \emph{nēmō, nēminis}
  \item \emph{Recall:} say the Latin for to live (somewhere), then the Latin for nothing, then say \emph{nēmō, nēminis} again
\end{itemize}

\begin{tcolorbox}[breakable,colback=teal!4,colframe=teal!35!black,title={Before you move on}]
What does nēmō, nēminis mean? (\textquotedblleft{}Nobody, no one\textquotedblright{}.) Say it once more.
\end{tcolorbox}
\section[nūllus, nūlla, nūllum]{nūllus, nūlla, nūllum — no, not any}

\label{lesson:LA-C161-nullus}



\begin{quote}
Before the new one: say the Latin for nothing, then the Latin for nobody.
\end{quote}

\subsection*{You'll want to know: nūllus, nūlla, nūllum}
\textbf{nūllus, nūlla, nūllum} — \textquotedblleft{}no, not any\textquotedblright{}. It agrees with its noun: \textbf{nūlla epistula} — \textquotedblleft{}no letter at all\textquotedblright{}. English \emph{null} and \emph{annul} come from it.

\begin{grammarlens}[title={What you've built}]
One more word for this chapter.
\end{grammarlens}

\subsection*{Your turn}
\begin{itemize}
\raggedright
  \item \emph{Say it:} \emph{nūllus, nūlla, nūllum}
  \item \emph{Say it:} \emph{nūllus, nūlla, nūllum}, once more
  \item \emph{Say it:} say \emph{nūllus, nūlla, nūllum}
  \item \emph{Recall:} say the Latin for nothing, then the Latin for nobody, then say \emph{nūllus, nūlla, nūllum} again
\end{itemize}

\begin{tcolorbox}[breakable,colback=teal!4,colframe=teal!35!black,title={Before you move on}]
What does nūllus, nūlla, nūllum mean? (\textquotedblleft{}No, not any\textquotedblright{}.) Say it once more.
\end{tcolorbox}
\section[uter, utra, utrum]{uter, utra, utrum — which (of two)}

\label{lesson:LA-C161-uter}



\begin{quote}
Before the new one: say the Latin for nobody, then the Latin for no.
\end{quote}

\subsection*{You'll want to know: uter, utra, utrum}
\textbf{uter, utra, utrum} — \textquotedblleft{}which (of two)\textquotedblright{}. A question that picks one of two: \textbf{Uter liber?} — \textquotedblleft{}Which of the two books?\textquotedblright{} \textbf{neuter}, \textquotedblleft{}neither of the two\textquotedblright{}, gave English \emph{neuter} and \emph{neutral}.

\begin{grammarlens}[title={What you've built}]
One more word for this chapter.
\end{grammarlens}

\subsection*{Your turn}
\begin{itemize}
\raggedright
  \item \emph{Say it:} \emph{uter, utra, utrum}
  \item \emph{Say it:} \emph{uter, utra, utrum}, once more
  \item \emph{Say it:} say \emph{uter, utra, utrum}
  \item \emph{Recall:} say the Latin for nobody, then the Latin for no, then say \emph{uter, utra, utrum} again
\end{itemize}

\begin{tcolorbox}[breakable,colback=teal!4,colframe=teal!35!black,title={Before you move on}]
What does uter, utra, utrum mean? (\textquotedblleft{}Which (of two)\textquotedblright{}.) Say it once more.
\end{tcolorbox}
\section[neque]{neque — and not, nor}

\label{lesson:LA-C161-neque}



\begin{quote}
Before the new one: say the Latin for no, then the Latin for which.
\end{quote}

\subsection*{You'll want to know: neque}
\textbf{neque} — \textquotedblleft{}and not, nor\textquotedblright{}. \textbf{neque ... neque} is \textquotedblleft{}neither ... nor\textquotedblright{}: \textbf{Neque aquam neque vīnum bibō.} — \textquotedblleft{}I drink neither water nor wine.\textquotedblright{} It is often shortened to \textbf{nec}.

\begin{grammarlens}[title={What you've built}]
One more word for this chapter.
\end{grammarlens}

\subsection*{Your turn}
\begin{itemize}
\raggedright
  \item \emph{Say it:} \emph{neque}
  \item \emph{Say it:} \emph{neque}, once more
  \item \emph{Say it:} say \emph{neque}
  \item \emph{Recall:} say the Latin for no, then the Latin for which, then say \emph{neque} again
\end{itemize}

\begin{tcolorbox}[breakable,colback=teal!4,colframe=teal!35!black,title={Before you move on}]
What does neque mean? (\textquotedblleft{}And not, nor\textquotedblright{}.) Say it once more.
\end{tcolorbox}
